\section{Discussion}\label{discussion}

Our study's first contribution to the science of spatiotemporal forecasting of opioid-related overdose deaths is highlighting the need for extensive comparisons to a robust suite of simple baselines. This lesson matches reports\citep{handClassifierTechnologyIllusion2006,rudinStopExplainingBlack2019} from across the sciences, especially efforts in health\citep{christodoulouSystematicReviewShows2019,nusinoviciLogisticRegressionWas2020} and the social sciences\citep{salganikMeasuringPredictabilityLife2020}, that suggest advanced modeling techniques may not substantially outperform simpler baselines on some difficult prediction tasks. Our findings are similar in both the large state of MA and the far denser Cook County. In each catchment area, across both intervention-aware and error-based metrics, we found that a historical average baseline performed competitively (within the uncertainty bounds of top-ranked statistical models). The key to success here is careful selection of the number of recent years in the look-back period, following standard best practices for hyperparameter tuning.\citep{probstTunabilityImportanceHyperparameters2018,raschkaModelEvaluationModel2020} If this simple baseline model yields such high performance, it raises questions about the rationale for adopting more complex counterparts that require specialized expertise. Many prior overdose forecasting studies\citep{allenTranslatingPredictiveAnalytics2023,bauerSmallAreaForecasting2023} completely omit such baselines, or often include only the poor performing ones such as the last-year\citep{marksIdentifyingCountiesRisk2021} model or a too-long historical average\citep{ertugrulCASTNetCommunityAttentiveSpatioTemporal2019}. For all future studies of opioid overdose forecasting, we recommend including historical averages with tuned look-back periods.

Our second contribution, developed in parallel to contemporary work\citep{marshallPreventingOverdoseUsing2022}, is a new metric -- percentage of best possible reach (\%BPR) -- which evaluates predictions based on their utility for informing decisions about where to intervene. In both Massachusetts and Cook County, we demonstrate that using \%BPR as an evaluation metric can lead to different model rankings and different recommendations of where to intervene than error-based metrics like RMSE, improving the total number of annual fatal overdose events that could be preemptively identified by 15 in IL and 18 in MA. This is an important finding, as we believe that intervention-aware metrics like \%BPR more closely reflect how public health agencies wish to use forecasting models to inform their intervention strategies.\citep{allenTranslatingPredictiveAnalytics2023}

Lastly, we emphasize that our study is designed to be reproducible and open to extensions by other researchers. We released the software for fitting all models and computing all metrics under a permissive open-source license (link in Introduction). We also released our cleaned version of the public-domain Cook County dataset as well as all preprocessing code. Historically, overdose forecasting studies have not often shared code and have focused on private bespoke datasets, often reasonably due to privacy issues around decedent data. Enabling diverse researchers to pursue a common prediction task, especially via the availability of a public dataset for evaluation, has been a key driver of progress in predictive modeling\citep{donoho50YearsData2017}.

\subsection{Limitations}\label{limitations}

This study has several limitations. First, our findings come from only two places (Massachusetts and Cook County), and may not be generalizable to other counties, states, or public health jurisdictions. Cook County is predominantly urban, while Massachusetts is a large state with substantial urban, rural and suburban areas. The spatiotemporal trends in opioid-related mortality could thus be dramatically different in these two locations, necessitating different model rankings and intervention strategies. Furthermore, we acknowledge that not all Cook County deaths are reported to the Medical Examiner. The Medical Examiner\textquotesingle s jurisdiction only covers specific fatalities for cause-of-death determination.

Second, there are limitations to our analysis of the proposed BPR metric. For simplicity, all results here assumed an intervention budget of K=100 census tracts. Different K values may lead to different method rankings. Our suggested BPR metric is intended for identifying where to intervene to relieve high overall burden. However, it does not directly prioritize the rate of change. Interventions aimed to reduce risk in communities that are at very high risk but do not already have a high burden may not be correctly identified using BPR.

Finally, other choices of covariates are possible. Our focus on a limited set of covariates, derived from the SVI of the American Community Survey, was an intentional choice to ensure the nationwide availability of these covariates. Certain jurisdictions may possess useful alternate data sources, such as emergency medical service (EMS) calls, insurance claims data, and measures for a mix of linked administrative datasets\citep{bharelUsingDataGuide2020}, necessitating additional covariate consideration for enhanced model performance.

\subsection{Conclusion}\label{conclusion}

In an effort to better predict future fatal opioid-related overdose spikes and inform future harm-reducing interventions, we compared overdose forecasting options. Our study reinforces the value of intervention-aware metrics like \%BPR in evaluating models for opioid overdose mortality forecasting. Our study also suggests that simple baselines like (weighted) historical averages should be included in future analyses, as more sophisticated and expensive-to-train models may not substantially outperform these baselines. As the opioid crisis continues to evolve, we hope our findings and our open-source resources enable improved model comparisons and better data-informed public health interventions that ultimately reduce the harm caused by overdose events.
